\section{Discussion and conclusion}\label{sec:discussion}
The analysis identifies several distinctions that an adaptive
trace-estimation policy must respect. Attempted sketch width is not accepted
rank. Accepted rank is not dominant-subspace capture. An ideal spectral tail
does not in general equal realized probe-specific residual energy. A useful
empirical comparison is not a finite-sample certificate, and safe acceptance
is not the same as a safe complete policy after information costs have been
paid.

The capture results explain two different failures. An invertible square
signal sketch can be poorly conditioned, amplifying a small nonzero tail.
A discrete signal sketch can also miss a direction exactly, even with
positive oversampling and a full-rank accepted basis. The first mechanism
helps explain tail-sensitive empirical mean curves; the second supplies an
exact local pilot-transcript limitation. Neither implies a universal
frequency of failure or a fixed oversampling amount that is always optimal.

The certification results also require two levels of interpretation.
Ordinary sample variance was useful under the preregistered empirical
decision gate. Charging for certification revealed a tradeoff between
aggregate mean-risk protection and harm on many ordinary paths. The explicit
confidence routes considered here were too conservative at the tested
sample sizes. This gap between empirical signal and available guarantees
is a research question, not a license to interpret an empirical guard as
a theorem.

\paragraph{Limitations.}
The conditional-risk bridges use finite frozen orientations and path
populations. Their bootstrap intervals do not quantify unseen rare-event
probabilities or orientation-universal performance. The Hessian and
pilot-gate experiments are exploratory. The structural-prior certificate
assumes valid prior information and does not establish that this information
can be obtained within the same budget. Our route-specific negative results
do not exclude other confidence constructions.

\paragraph{Open problem.}
A focused continuation is to derive an observable, budget-feasible
one-sided upper confidence bound for the common-probe difference
\[
\Delta_R=\frac{\sigma_a^2}{\ell_a}-\frac{\sigma_0^2}{\ell_0}.
\]
Such a result should exploit paired information when justified, charge all
committed construction and certification queries, preserve the freshness of
final residual probes, and specify a feasible fallback. A fresh-query
capture test or a genuine small-ball analysis is an alternative direction.
No new allocator is proposed or implemented in this manuscript.

In summary, the project changes the emphasis from predicting a favorable
split from visible spectral features to justifying a decision using the
residual risk that the realized estimator actually incurs.
